%=============================================================================
\section{实战调试与调参手册}
%=============================================================================

这是你比赛当天真正会翻开的一章。前面那些章节教你怎么设计控制器；这一章教你——控制器炸了怎么办。这里的调试技巧不是纸上谈兵，它们来自几百个小时的机器人翻车、电机烧毁和赛前最后一秒的抢修。

你的机器人刚刚开始剧烈抖动。离比赛还有二十分钟。你没时间重新推导传递函数，也没空回头看 Bode 图那章。你需要一套系统化的流程，让你从"好像有大问题"尽快走到"修好了"。这就是本章要给你的东西。

%-----------------------------------------------------------------------------
\subsection{调试心态}
%-----------------------------------------------------------------------------

在讲具体症状和修法之前，先把三条铁律刻进脑子里。违反其中任何一条，你都会白白浪费好几个小时。

\textbf{铁律一：机器人永远没错。}它振荡了，一定有原因。它漂移了，一定有原因。物理不会骗你；是你对物理的建模出了问题。不要说"电机表现很奇怪"，要说"电机在做 X，原因是 Y，我需要改 Z"。机器人严格执行你告诉它的指令——问题在于你告诉它的和你想让它做的是不是同一件事。

\textbf{铁律二：一次只改一样东西。}如果你同时改了增益又改了硬件，你永远不会知道到底是哪个修好了问题。如果你同时改了 $K_p$ 和 $K_d$，振荡消失了，是 $K_p$ 太高还是 $K_d$ 太低？你不知道。规范的调试是：改一个参数、测试、观察、记录。这看起来慢，但比另一种做法快得多——在参数空间里瞎撞。

\textbf{铁律三：记录一切。}"我觉得电机好像在抖"没有任何价值。"速度在 42~Hz 处振荡，幅值 $\pm 200$~RPM"才是可操作信息。好的日志让你事后分析问题、比较改前改后的行为、跟队友精确沟通。如果你的 MCU 有 UART 或 USB 日志功能，用它。如果你有逻辑分析仪，用它。实在没有，就在 ISR 里翻转一个 GPIO 引脚，拿示波器看时序。

%-----------------------------------------------------------------------------
\subsection{诊断流程图}
%-----------------------------------------------------------------------------

出了问题，第一件事是识别\textit{症状}。下面列出最常见的几种症状和对应的诊断流程。

%- - - - - - - - - - - - - - - - - - - - - - - - - - - - - - - - - - - - - - -
\subsubsection{症状：振荡（Oscillation）}
%- - - - - - - - - - - - - - - - - - - - - - - - - - - - - - - - - - - - - - -

你的机器人在抖、在震、在嗡嗡响。这是比赛当天出现频率最高的问题。

\begin{enumerate}
    \item \textbf{是机械问题吗？}断开控制器，给一个恒定的开环占空比，用手转转电机。机构本身会不会振动？如果会，问题不在你的控制器——去拧螺丝、查轴承、消间隙。

    \item \textbf{是传感器噪声吗？}电机静止时记录原始传感器数据。噪声大吗？如果编码器信号跳了好几个 count，或者 IMU 读数在来回蹦，那这些噪声正在被你的微分项放大。给传感器信号加个低通滤波器（第2章）。

    \item \textbf{是 $K_p$ 太高了吗？}把 $K_p$ 减半。振荡停了还是明显减弱了？如果是，说明你已经超出了稳定裕度。继续降 $K_p$ 直到稳定，然后小心地增大 $K_d$ 来恢复性能。

    \item \textbf{是采样率太低了吗？}在控制 ISR 的开头和结尾翻转一个 GPIO 引脚，用示波器量脉宽。实际循环周期和你设计的 $T_s$ 一致吗？如果循环跑得比你以为的慢，你的有效增益就跟设计值不一样了。

    \item \textbf{是结构共振吗？}对振荡信号做 FFT（或者至少在示波器上数周期估个频率）。振荡集中在某个与 PID 带宽无关的尖锐频率上吗？如果是，你在激励一个机械共振。在那个频率上加陷波滤波器（notch filter，第2章），或者降低 $K_d$——微分项最容易激发共振。
\end{enumerate}

%- - - - - - - - - - - - - - - - - - - - - - - - - - - - - - - - - - - - - - -
\subsubsection{症状：超调（Overshoot）}
%- - - - - - - - - - - - - - - - - - - - - - - - - - - - - - - - - - - - - - -

输出越过设定值才回来。适度的超调（5--10\%）是正常的，通常可以接受。过大的超调（>20\%）就说明有问题了。

\begin{enumerate}
    \item \textbf{是 $K_d$ 太低了吗？}微分项就是你的阻尼。增大 $K_d$——超调有没有减小？如果你根本没有微分项，那你跑的是 PI 控制器，面对阶跃输入天然就会超调。

    \item \textbf{是 $K_i$ 太高了吗？}积分项在瞬态过程中持续累积，累积值会把输出推过设定值。降低 $K_i$，或者加积分分离（integral separation，第4章）——误差大的时候关掉积分项，只在接近设定值时才启用。

    \item \textbf{是执行器饱和了吗？}记录 clamp 之前的原始控制器输出。如果它经常打到最大值，说明控制器在要求执行器给出超过能力的输出。饱和期间系统实际上是开环的，而积分项还在往上加。加 anti-windup（第4章）。

    \item \textbf{是设定值阶跃太激进了吗？}从 0 跳到 5000~RPM 意味着在 $t=0$ 时要求无穷大加速度。没有物理系统能做到。给参考信号加个设定值滤波器（setpoint filter）或者软启动斜坡（soft-start ramp，第2章），限制参考信号的变化率。
\end{enumerate}

%- - - - - - - - - - - - - - - - - - - - - - - - - - - - - - - - - - - - - - -
\subsubsection{症状：稳态误差 / 漂移（Steady-State Error / Drift）}
%- - - - - - - - - - - - - - - - - - - - - - - - - - - - - - - - - - - - - - -

输出稳住了，但不在设定值上。有一个持续的偏差。

\begin{enumerate}
    \item \textbf{有积分项吗？}如果 $K_i = 0$，纯 PD 控制器在存在恒定扰动（重力、摩擦、负载）时一定会有稳态误差。加个积分项。

    \item \textbf{传感器有偏置吗？}检查传感器标定（第2b章）。一个有 0.5\textdegree{} 偏置的 IMU，无论你的控制器多好，都会造成 0.5\textdegree{} 的稳态误差。

    \item \textbf{有机械偏移吗？}重力、皮带张力、弹簧预紧力——任何你的控制器模型没有考虑到的恒定力，都会表现为稳态误差。要么加前馈补偿（feedforward compensation），要么让积分项来吸收它（但这样更慢）。

    \item \textbf{积分饱和（integral windup）反而阻止了积分器做它该做的事？}说来矛盾：如果你的 anti-windup 太激进，可能在积分累积到足够消除误差之前就把它截住了。检查你的 anti-windup 限幅（第4章）是否大到能应对预期的稳态负载。
\end{enumerate}

%- - - - - - - - - - - - - - - - - - - - - - - - - - - - - - - - - - - - - - -
\subsubsection{症状：没反应（No Response）}
%- - - - - - - - - - - - - - - - - - - - - - - - - - - - - - - - - - - - - - -

你给电机发了运动指令。什么都没发生。这要么是最好查的问题，要么是最难查的。

\begin{enumerate}
    \item \textbf{电机有电吗？}拿万用表量电机端子电压。这听起来很蠢，但保险丝烧了、接插件松了、电池没电了——这类问题在"我的控制器不工作"的投诉里占了惊人的比例。

    \item \textbf{通信链路通了吗？}对于 CAN 总线无刷电调（C610/C620），检查：CAN 总线初始化了吗？发送的 CAN ID 和电调期望的一致吗？电调状态灯正常闪烁吗？用逻辑分析仪或 CAN 调试器确认帧确实在总线上。对于 PWM 驱动的电调，检查使能引脚并用示波器验证 PWM 信号。

    \item \textbf{指令极性对吗？}正电流指令应该产生你期望方向的正力矩。如果电机一上电就飞车，说明符号约定反了。C610/C620 的 16 位电流指令是有符号数——正负对应相反的力矩方向。

    \item \textbf{控制循环真的在跑吗？}在控制循环 ISR 里翻转一个 GPIO 引脚。用示波器看它是不是按预期频率翻转。如果根本没翻转，说明你的定时器中断配置不对。

    \item \textbf{符号约定对吗？}这是最危险的软件 bug。如果正误差产生了\textit{负}输出，你的系统就是正反馈回路。它不会温和地振荡——它会瞬间冲到极限。检查：正的设定值变化是否产生正的控制输出？
\end{enumerate}

%-----------------------------------------------------------------------------
\subsection{机械？电气？软件？还是控制？}
%-----------------------------------------------------------------------------

电机还在抖，你已经把 $K_p$ 减到了 0.1——抖动完全没变化。这时候该怀疑的就不是控制参数了。很多调试失败的根本原因是：学生在调 PID 增益，而真正的问题是一颗松了的顶丝、一个电压不够的电源、或者代码里的一个符号错误。在动增益之前，先给问题分个类。

%- - - - - - - - - - - - - - - - - - - - - - - - - - - - - - - - - - - - - - -
\subsubsection{隔离测试（The Isolation Test）}
%- - - - - - - - - - - - - - - - - - - - - - - - - - - - - - - - - - - - - - -

这是最有价值的单项调试技术。两分钟搞定，能省你好几个小时。

\textbf{第一步：}断开控制器。给一个恒定的开环力矩（比如通过 CAN 发送固定电流指令给 C610/C620，或者对简单驱动器给一个恒定 PWM 占空比）。观察电机。
\begin{itemize}
    \item 如果电机转得平稳又匀速：问题出在控制器或软件上。去查软件/控制相关的项目。
    \item 如果电机抖动、卡死、或者行为异常：问题出在机械或电气上。先别动代码。
\end{itemize}

\textbf{第二步：}查机械。把电机从机构上完全拆下来。用手转输出轴。卡不卡？间隙大不大？手感粗糙（轴承坏了）还是一顿一顿的（皮带张力、齿轮啮合）？

\textbf{第三步：}查电气。电机带载运行时量一下电源电压。压降大不大（超过 10\%）？如果多个轴同时运动时机器人表现更差，几乎可以肯定是电源问题。

%- - - - - - - - - - - - - - - - - - - - - - - - - - - - - - - - - - - - - - -
\subsubsection{常见电气问题}
%- - - - - - - - - - - - - - - - - - - - - - - - - - - - - - - - - - - - - - -

我们来看三类最常见的电气陷阱——它们的共同特征是：症状像控制问题，但调增益完全没用。

\begin{itemize}
    \item \textbf{负载下电源电压跌落。}电池或电源供不上足够的电流。多个电机同时吸电时电压下降，所有控制器都丧失控制力。症状是：单个电机跑没问题，所有电机一起跑就炸。修法：用容量更大的电池、在电机驱动附近加大容量电容、或者把电机启动时间错开。

    \item \textbf{地环路（Ground loop）。}如果你的逻辑板和功率板共用地线，电机电流流过共地路径会产生压降，表现为传感器噪声。噪声与电机电流成正比——电机越使劲噪声越大。修法：使用分离的地平面、星形接地、或者光耦隔离传感器接口。

    \item \textbf{PWM 开关的 EMI。}PWM 开关的快沿（dV/dt 可以超过 1~V/ns）会产生电磁干扰，干扰附近的传感器信号。模拟传感器（电流采样、电位器）和慢速数字总线（I2C）尤其容易中招。修法：传感器线远离电机线、使用屏蔽线缆、在传感器输入端加滤波电容、对噪声敏感的场合优选 SPI 而非 I2C。

    \item \textbf{连接器接触电阻。}腐蚀或松动的连接器可以在功率通路上引入 0.1--1~$\Omega$ 的电阻。10~A 的电机电流下就是 1--10~V 的压降——足以饿死电机或让 MCU 掉电。症状是间歇性的：实验台上好好的，比赛负载下就出问题。电机运行时晃晃每个电源连接器，看有没有毛刺。修法：清洁连接器、使用压接良好的 XT60/XT30 连接器、在每个电机驱动附近加一个大电容（100--470~$\mu$F）来缓冲瞬态电流需求。
\end{itemize}

%- - - - - - - - - - - - - - - - - - - - - - - - - - - - - - - - - - - - - - -
\subsubsection{常见软件问题}
%- - - - - - - - - - - - - - - - - - - - - - - - - - - - - - - - - - - - - - -

\begin{itemize}
    \item \textbf{符号约定错误（Sign convention error）。}如果正电压产生的旋转方向是让误差\textit{增大}而不是减小，你就构成了正反馈。系统不会振荡——它会直接跑飞。这是最常见也最危险的软件 bug。永远验证：正误差 $\to$ 正输出 $\to$ 误差减小。

    \item \textbf{变量溢出（Variable overflow）。}16 位编码器计数器从 65535 翻回 0。如果你计算速度时用 \texttt{count\_now - count\_prev} 却没处理溢出，就会得到一个巨大的尖峰。类似地，积分项如果用 16 位整数累加也会溢出。所有 PID 计算都用 32 位（或浮点）算术。

    \item \textbf{时序错误（Timing error）。}你的控制循环跑在错误的频率上，因为定时器预分频器配错了。你按 1~kHz 设计增益，但循环实际跑在 500~Hz。你的有效 $K_d$ 是设计值的两倍（因为 $K_d$ 与采样率有关），系统会激进得多。

    \item \textbf{数据竞争（Data race）。}ISR 更新一个共享变量（编码器计数、传感器读数）时，主循环正在读它。如果没有 \texttt{volatile} 关键字（C/C++）或者正确的互斥锁，编译器可能优化掉读操作，或者你读到一个写了一半的值。共享变量永远声明为 \texttt{volatile}，读取多字节共享值时短暂关中断。安全读取的示例：

\begin{verbatim}
volatile int32_t encoder_count;  // written in ISR

int32_t read_encoder(void) {
    __disable_irq();
    int32_t val = encoder_count;
    __enable_irq();
    return val;
}
\end{verbatim}

    \item \textbf{未初始化变量（Uninitialized variable）。}如果你忘了把积分累加器初始化为零，它的初始值就是内存里的垃圾。这可能导致启动时的大瞬态——电机猛抽一下然后稳住。初始化函数里永远把所有 PID 状态变量清零。

    \item \textbf{单位或缩放错误（Wrong units or scaling）。}传感器返回的是 count，但控制器期望的是弧度（radian）。PWM 寄存器接受 0--999，但 PID 输出的是 0--100。单位不匹配会让增益偏差好几个数量级。在每个接口边界都用注释标明单位。
\end{itemize}

%-----------------------------------------------------------------------------
\subsection{示波器与逻辑分析仪技巧}
%-----------------------------------------------------------------------------

你的眼睛和耳朵是很差的测量仪器。"好像在抖"不能告诉你任何定量信息。合适的工具能让不可见的问题变得可见。

\textbf{指令验证。}对于 CAN 总线系统（C610/C620），用 CAN 调试器或逻辑分析仪解码总线上的帧。验证：(1) MCU 按预期速率发送（通常 1~kHz），(2) CAN ID 与电调期望的一致，(3) 电流指令字节包含你期望的值。对于 PWM 驱动器，用示波器验证占空比是否与软件指令匹配。

\textbf{控制循环时序。}在 ISR 开头翻转一个 GPIO 引脚，结尾清除。在示波器上，脉宽就是你的 ISR 执行时间，周期就是循环周期。两者都应该与设计一致。如果脉宽接近周期，说明你的 ISR 太慢——它会错过 deadline，控制质量下降。代码如下：

\begin{verbatim}
void TIM_IRQHandler(void) {
    HAL_GPIO_WritePin(DEBUG_GPIO, GPIO_PIN_SET);   // ISR start

    // ... your control code here ...

    HAL_GPIO_WritePin(DEBUG_GPIO, GPIO_PIN_RESET); // ISR end
}
\end{verbatim}

\textbf{传感器信号完整性。}用逻辑分析仪探 I2C/SPI 总线流量。验证数据字节与你软件读到的值一致。如果总线上数据正确但你的变量值不对，问题在解析代码里。如果总线本身有毛刺，问题在电气层面（噪声、时钟速度太高、I2C 缺上拉电阻）。

\textbf{记录 PID 内部变量。}卡住的时候，记录\textit{所有} PID 内部变量，不要只看输出。一条有用的日志行应该打印：时间戳、设定值、测量值、误差、P 项、I 项、D 项、原始输出和截幅后输出。下面是一个最小 UART 日志示例：

\begin{verbatim}
// Call once per control cycle (or every N cycles to reduce bandwidth)
void log_pid(uint32_t tick, float sp, float meas,
             float p, float i, float d, float out_raw, float out) {
    printf("%lu,%.1f,%.1f,%.2f,%.2f,%.2f,%.1f,%.1f\r\n",
           tick, sp, meas, p, i, d, out_raw, out);
}
\end{verbatim}

测试完用 Python（甚至 Excel）画图。P/I/D 的分项拆解能立刻揭示哪一项在捣乱。如果 I 项很大而 P 项很小，说明有 windup。如果 D 项噪声很大，说明需要滤波。这一项技术能帮你省的时间，比其他任何技术都多。

\textbf{电流测量。}如果你有采样电阻（current-sense resistor）或霍尔效应传感器（Hall-effect sensor），记录实际电机电流极其有价值。电流告诉你电机\textit{实际在做什么}，而不是你\textit{命令它做什么}。命令电流和实际电流之间的不匹配，能揭示驱动器问题、接线问题或电机故障。

\textbf{小贴士：}一个 15 美元的逻辑分析仪（Saleae 兼容版）是你能为比赛机器人买到的最超值调试工具。它能解码 I2C、SPI、UART、CAN 和 PWM。配合 PulseView（免费软件），它能替代几百美元的设备完成 90\% 的调试任务。

%-----------------------------------------------------------------------------
\subsection{真实硬件上的 PID 调参：完整实例}
%-----------------------------------------------------------------------------

理论再好，你也需要一套能在工作台上一步一步跟着走的具体流程。

%- - - - - - - - - - - - - - - - - - - - - - - - - - - - - - - - - - - - - - -
\subsubsection{系统：DJI 3508 电机速度控制}
%- - - - - - - - - - - - - - - - - - - - - - - - - - - - - - - - - - - - - - -

DJI 3508 电机搭配 C620 无刷电调在 RoboMaster 比赛中随处可见。C620 是 CAN 总线设备：你的 MCU 通过 CAN 发送 16 位有符号电流（力矩）指令，C620 内部处理电流环、FOC 换相和编码器读取，并通过 CAN 回传转子角度、转速和实际电流。相关参数如下：
\begin{itemize}
    \item 转子惯量：约 $J = 0.0005$~kg$\cdot$m$^2$
    \item 力矩常数：$K_t \approx 0.3$~N$\cdot$m/A
    \item 绕组电阻：$R \approx 0.5~\Omega$
    \item CAN 协议：一帧 TX 携带 4 个电机的电流指令；每个电机一帧 RX 返回角度、转速和电流反馈。
    \item 电流环：由 C620 内部以 $\sim$10~kHz 处理。我们把它视为快速内环，在 MCU 上只设计外层速度环。
    \item 目标：速度跟踪，超调 $<5\%$，调节时间 $<200$~ms。
\end{itemize}

%- - - - - - - - - - - - - - - - - - - - - - - - - - - - - - - - - - - - - - -
\subsubsection{逐步调参流程}
%- - - - - - - - - - - - - - - - - - - - - - - - - - - - - - - - - - - - - - -

\textbf{第一步：找临界增益。}设 $K_i = 0$，$K_d = 0$。发一个 500~RPM 的阶跃指令。从 0 开始慢慢增大 $K_p$。到某个值时，速度会出现持续的、等幅振荡。这就是\textit{临界增益（critical gain）}。记录 $K_{p,\text{osc}} \approx 15$。然后设 $K_p = 0.5 \times K_{p,\text{osc}} = 7.5$。系统应该稳定但响应偏慢。

\textbf{第二步：加微分项。}从 $K_d = 0.01$ 开始，逐渐增大。微分项提供阻尼，减小超调并改善调节时间。持续增大直到超调低于 5\%。典型结果：$K_d \approx 0.05$。如果增大 $K_d$ 引起高频振动，说明你在放大传感器噪声——给微分通道加个低通滤波器。

\textbf{第三步：慢慢加积分项。}从 $K_i = 0.5$ 开始。发同样的 500~RPM 阶跃。检查：稳态误差消失了吗？增大 $K_i$ 直到稳态误差可以忽略。典型结果：$K_i \approx 2.0$。如果增大 $K_i$ 导致超调回来了，说明加过了——退回去接受一点稳态误差，或者加积分分离（integral separation）。

\textbf{第四步：扰动测试。}电机稳定运行在 500~RPM 时，用手短暂加载输出轴（或者拿块橡皮按住输出端）。速度应该掉一下然后在 200--300~ms 内恢复。如果恢复太慢，稍微增大 $K_i$。如果恢复过程中振荡了，降低 $K_p$ 或增大 $K_d$。

\textbf{第五步：全工况测试。}在 100、500、2000 和 5000~RPM 分别重复阶跃响应。电机的行为随转速变化（反电动势、摩擦力）。如果控制器在 500~RPM 正常但在 5000~RPM 振荡，你可能需要增益调度（gain scheduling）或者一组更保守的、能覆盖全范围的增益。

调完这套流程，你应该得到一个速度响应：快速上升到设定值、超调不超过 5\%、200~ms 内稳定。阶跃响应曲线应该是一条平滑的上升加极少的振铃——不是慢吞吞地爬上去（欠调），也不是疯狂振荡（过调）。

\textbf{每一步的响应曲线长什么样：}
\begin{itemize}
    \item 第一步之后（只有 $K_p$）：上升慢，约 30\% 超调，可见振铃（稳定前振荡 2--3 次）。稳态误差大约 5--10\%。
    \item 第二步之后（$K_p + K_d$）：上升更快，超调降到 $<5\%$，约 150~ms 稳定。但仍有持续的稳态误差。
    \item 第三步之后（$K_p + K_i + K_d$）：上升速度和超调不变，但误差在 200~ms 内收敛到零。积分项平滑地补上了最后几个百分点。
    \item 扰动测试：加载时速度掉了大约 50~RPM，200--300~ms 内完全恢复。
\end{itemize}

%- - - - - - - - - - - - - - - - - - - - - - - - - - - - - - - - - - - - - - -
\subsubsection{Ziegler-Nichols 失效的场景}
%- - - - - - - - - - - - - - - - - - - - - - - - - - - - - - - - - - - - - - -

Ziegler-Nichols 极限增益法（也就是上面第一步的基础）是一个有用的起点，但在以下几种重要场景中会失败：

\begin{itemize}
    \item \textbf{有大时延的系统。}如果传输延迟与系统时间常数可比，Z-N 给出的增益会过于激进。得到的控制器几乎没有相位裕度。

    \item \textbf{有非线性摩擦的系统。}如果静摩擦力（stiction）显著，振荡测试在不同幅值下会给出不一致的结果。低幅值时电机粘住不动，高幅值时转得平滑。"临界增益"取决于测试幅值，这就失去了意义。

    \item \textbf{振荡有危险的系统。}对于平衡机器人或重型机械臂，在临界增益处诱发持续振荡可能造成机械损坏或人身伤害。
\end{itemize}

在以上所有情况下，用前面那套手动逐步调参流程。从很低的增益开始，逐步增大，依靠阶跃响应测量而不是振荡边界。这更麻烦，但安全得多。

%-----------------------------------------------------------------------------
\subsection{常见故障模式及其特征}
%-----------------------------------------------------------------------------

学会在数据日志里识别这些模式。每种故障都有独特的"指纹"——一旦你见过一次，以后一眼就能认出来。

%- - - - - - - - - - - - - - - - - - - - - - - - - - - - - - - - - - - - - - -
\subsubsection{积分饱和的实际表现（Integral Windup in Practice）}
%- - - - - - - - - - - - - - - - - - - - - - - - - - - - - - - - - - - - - - -

\textbf{特征：}执行器饱和之后，一个方向的响应变得迟钝，等到误差反向时出现猛烈超调。

\textbf{时间序列长什么样：}输出指令打到最大限幅并持续一段时间。在此期间积分项持续累积。当误差终于变号时，积分项大到需要很久才能"泄放"，导致输出在本该反向之后很久仍然维持在最大值。结果是反方向出现巨大超调。

\textbf{修法：}实现 clamping 或 back-calculation anti-windup（第4章）。Clamping 方式：输出饱和时停止积分累积。Back-calculation 方式：将饱和输出与未饱和输出的差值反馈回来，加速积分器泄放。

%- - - - - - - - - - - - - - - - - - - - - - - - - - - - - - - - - - - - - - -
\subsubsection{执行器饱和（Actuator Saturation）}
%- - - - - - - - - - - - - - - - - - - - - - - - - - - - - - - - - - - - - - -

\textbf{特征：}低幅值时实际响应与仿真吻合，高幅值时出现偏差。设定值越高性能越差。

\textbf{诊断：}记录 clamp 之前的原始控制器输出。如果它经常超出执行器限幅（比如 PID 输出 $\pm 20000$，但电机驱动只接受 $\pm 16384$），控制器在要求执行器做超出能力的事。饱和期间系统实际上是开环的。

\textbf{修法：}降低增益使输出保持在线性范围内、加带 anti-windup 的输出限幅、或者接受在极端工况下性能变差。有时候答案就是"你的执行器对于这个任务来说太弱了"——再怎么调参也改变不了物理。

%- - - - - - - - - - - - - - - - - - - - - - - - - - - - - - - - - - - - - - -
\subsubsection{传感器混叠（Sensor Aliasing）}
%- - - - - - - - - - - - - - - - - - - - - - - - - - - - - - - - - - - - - - -

\textbf{特征：}出现一个意料之外频率的幽灵振荡，跟 PID 带宽和机械共振都不搭边。

\textbf{原因：}传感器采样率相对于信号带宽太低。一个高频信号（比如 800~Hz 的电机振动）以 1~kHz 采样就会混叠到 200~Hz（Nyquist 定理，第2章）。控制器看到了一个物理上根本不存在的 200~Hz 扰动，然后试图去抑制它。

\textbf{修法：}把传感器采样率提高到最高信号频率的至少 $2\times$。或者在 ADC 之前加一个模拟抗混叠滤波器（analog anti-aliasing filter，一个简单的 RC 低通就行），在采样前去掉高频成分。

%- - - - - - - - - - - - - - - - - - - - - - - - - - - - - - - - - - - - - - -
\subsubsection{量化极限环（Quantization Limit Cycling）}
%- - - - - - - - - - - - - - - - - - - - - - - - - - - - - - - - - - - - - - -

\textbf{特征：}围绕设定值的小幅持续振荡，永远不衰减。幅值通常是 1--2 个编码器 count。

\textbf{原因：}编码器分辨率相对于你要求的定位精度太粗。控制器看到"我在设定值上方 1 个 count"，就发一个修正指令。电机移到下方 1 个 count。现在低了 1 个 count，又修正回去。这种来回交替就是量化引起的\textit{极限环（limit cycle）}。

\textbf{修法：}在设定值附近加一个小的死区（deadband）——忽略低于量化阈值的误差——或者换一个分辨率更高的编码器。死区宽度大约设为 1--2 个编码器 count。

%- - - - - - - - - - - - - - - - - - - - - - - - - - - - - - - - - - - - - - -
\subsubsection{通信延迟（Communication Latency）}
%- - - - - - - - - - - - - - - - - - - - - - - - - - - - - - - - - - - - - - -

\textbf{特征：}CAN 总线负载增大或更多设备共享通信总线时，不稳定性加剧。

\textbf{原因：}传感器读数的延迟不确定，意味着控制器在用过时的数据做决策。如果延迟与循环周期可比，就会增加显著的相位滞后。相位裕度降低，系统趋向不稳定。总线负载重时，消息排队、延迟增大，系统就失稳了。

\textbf{修法：}降低控制器带宽（降低 $K_p$ 和 $K_d$），使环路对延迟更鲁棒。或者用 Smith predictor（第4章）做显式延迟补偿——建模已知延迟并在控制律中补偿它。


%- - - - - - - - - - - - - - - - - - - - - - - - - - - - - - - - - - - - - - -

\begin{mdframed}[linewidth=1pt, innertopmargin=10pt, innerbottommargin=10pt]
\textbf{案例：昨天还好好的平衡车}

\vspace{0.3em}
\textbf{场景：}你的平衡车在实验室光滑地面上完美运行。到了比赛场地，在地毯上 3 秒内就倒了。

\textbf{根因：}地毯的摩擦力比光滑地面大得多。轮子的抓地力不同了，有效阻尼变了，而且——最关键的——积分项的累积方式变了，因为摩擦力产生了更大的恒定扰动需要积分器去对抗。在光滑地面上积分器几乎不激活。在地毯上它在 2 秒内就饱和了。

\textbf{诊断：}你在两种地面上都记录了积分项。光滑地面上接近零。地毯上 2 秒内冲到饱和限幅并停在那里。Anti-windup 介入了，但那时机器人已经倾斜超过了可恢复角度。

\textbf{修法（2 分钟版）：}$K_d$ 加 50\%（地毯的物理阻尼部分替代了微分项的作用，但更高的摩擦力也造成了滞后，$K_d$ 有助于应对）。$K_i$ 减 40\% 以避免过早饱和。Anti-windup 限幅稍微放宽一点。

\textbf{修法（正经版）：}加一个摩擦力前馈项（friction feedforward），根据轮子速度的符号估计摩擦力并加入补偿力矩。这样就完全把负担从积分项上卸下来了。

\textbf{教训：}比赛前一定要在实际比赛场地上测试。而且随时准备好一套"高摩擦"参数组，可以一键切换。
\end{mdframed}

\vspace{1em}

\begin{mdframed}[linewidth=1pt, innertopmargin=10pt, innerbottommargin=10pt]
\textbf{案例：23~Hz 共振的云台}

\vspace{0.3em}
\textbf{场景：}你的云台在 23~Hz 抖。你很困惑，因为 PID 带宽才 10~Hz。没有任何 PID 参数跟 23~Hz 沾边。

\textbf{根因：}云台框架有一个 23~Hz 的结构共振。你的 PID 里的微分项在远超 PID 带宽的频率上仍然有增益（记住，微分放大高频）。在 23~Hz 处，微分项的增益足以激励共振，而共振通过 IMU 反馈回来，形成自激振荡。

\textbf{诊断：}对电机电流指令做 FFT。在 23~Hz 处出现一个尖峰。把 $K_d$ 减半——振动幅值按比例下降。这证实了微分项就是激励源。

\textbf{修法（快速版）：}降低 $K_d$，接受稍慢的瞬态响应。这是比赛日的修法。

\textbf{修法（正经版）：}在反馈通道里加一个 23~Hz 的陷波滤波器（notch filter，第2章）。陷波器把共振频率从 PID 看到的信号中移除，这样微分项就激励不了它了。频谱的其余部分不受影响，瞬态性能得以保留。设计陷波器的 Q 因子在 5--10 之间（窄到不影响邻近频率，宽到即使共振频率随温度或负载略有偏移也能兜住）。

\textbf{教训：}PID 带宽不是唯一重要的频率。微分项在带宽\textit{以上}的所有频率都有增益。任何在微分有效范围内的机械共振都是潜在受害者。设计云台时，永远先表征结构共振，\textit{然后}再调控制器。
\end{mdframed}

\vspace{1em}

%-----------------------------------------------------------------------------
\subsection{打造你的调试工具箱}
%-----------------------------------------------------------------------------

一个赛季下来，你应该积累起一套调试工具和习惯。下面是一个有经验的战队的工具箱长什么样：

\textbf{硬件工具：}
\begin{itemize}
    \item 万用表（查电压、电流、通断）
    \item 示波器（哪怕是像 Hantek 6022BE 这样的便宜 USB 示波器，也够看 PWM 和时序了）
    \item 逻辑分析仪（Saleae 兼容版 + PulseView，解码 I2C/SPI/CAN/UART）
    \item 备用电池、连接器和线材
    \item 一套已知正常的电机和驱动器（用于隔离测试）
\end{itemize}

\textbf{软件工具：}
\begin{itemize}
    \item 串口绘图工具（SerialPlot、PlotJuggler，或者一个简单的 Python + matplotlib 脚本）
    \item 参数调节界面——GUI 也好，简单的串口命令界面也好，只要能让你不重新刷固件就改 $K_p$、$K_i$、$K_d$
    \item 一个能以循环频率 dump 所有 PID 内部变量的日志框架
    \item 一个测试脚本，能跑标准阶跃响应并自动算出超调、调节时间和稳态误差
\end{itemize}

\textbf{软件参数在线调节示例。}一个最简的串口命令解析器，让你在运行时改增益：

\begin{verbatim}
// Parse "Kp=7.5" or "Ki=2.0" from UART input
void parse_gain_command(const char *cmd) {
    float val;
    if (sscanf(cmd, "Kp=%f", &val) == 1) pid.Kp = val;
    if (sscanf(cmd, "Ki=%f", &val) == 1) pid.Ki = val;
    if (sscanf(cmd, "Kd=%f", &val) == 1) pid.Kd = val;
}
\end{verbatim}

这省掉了"刷固件-重启-再测"的循环，每次迭代轻松吃掉 2 分钟。调 50 次就是一个多小时。

\textbf{调试日志。}维护一个简单的文本文件（或笔记本、或共享文档），记录每次调试：
\begin{itemize}
    \item 日期和时间
    \item 观察到的症状
    \item 测试的假设
    \item 改了什么
    \item 结果
\end{itemize}

这听起来很烦，但它能防止最常见的"元调试"问题："等等，我们是不是已经试过了？当时怎么样来着？"三个队友在比赛前一晚 11 点同时改东西的时候，一份日志就是系统性进步和原地打转的分界线。

\vspace{0.5em}

\textbf{最后一句话。}调试不代表你做错了什么。每个真实系统都需要调试。好工程师和焦头烂额的学生之间的区别，不在于 bug 的有无——在于找到并修复 bug 的速度。本章的技术不会消灭问题，但会大幅缩短从"好像有什么不对"到"我知道哪里不对以及怎么修"之间的时间。这个速度，就是赢比赛的关键。
